\chapter{Sound, Quantum Oscillators, and the Second Law}

\lectureoverview{The second law is the destination, but two concrete
calculations prepare the way.  The speed of sound tests how thermodynamic
derivatives become observable motion.  The harmonic oscillator then exposes
the failure of classical equipartition for stiff degrees of freedom and shows
how quantum mechanics activates modes one by one.  We finally return to the
apparent conflict between reversible microscopic dynamics and increasing
entropy, using phase-space flow, coarse-graining, chaos, and recurrence to
state the second law as an overwhelmingly probable law rather than an
exceptionless mechanical prohibition.}

\section{A Promise About the Second Law}

The central problem is easy to state and difficult to explain.  Newtonian
equations are reversible: if a microscopic history is allowed, the
time-reversed history is allowed as well.  The second law nevertheless says
that macroscopic processes have a preferred direction, characterized by
entropy that ordinarily increases.  We shall come to that tension, but not
immediately.  It is better first to see statistical mechanics produce
concrete information and then encounter, within an elementary calculation,
the historical puzzle that forced classical physics to give way to quantum
physics.

The first example is numerical and deliberately rough: estimate the speed of
sound in a dilute gas.  The second is exact within its assumptions: calculate
the mean energy of a harmonic oscillator in a heat bath.  The oscillator will
lead from equipartition to quantized energy levels, from one mode to infinitely
many modes, and finally back to the question of what it means for a
many-particle system to forget its initial state.

\section{The Molecular Scale of the Speed of Sound}

Imagine a small overdensity in a dilute gas.  It spreads because molecules
leave the crowded region, so its propagation speed should be comparable with
the molecular speed.  Equipartition gives the kinetic energy of one molecule,
\begin{equation}
 \frac{3}{2}k_B T
 =\frac{1}{2}m\expect{v^2},
 \qquad
 \expect{v^2}=\frac{3k_B T}{m}.
 \label{eq:sm7-molecular-speed}
\end{equation}
Here temperature is measured from absolute zero and \(k_B\) has been restored
because we want laboratory units.  The argument is only a scale estimate:
the molecular root-mean-square speed is not itself the speed of a collective
sound wave.

A more systematic derivation considers a small element of gas.  Unequal
pressures on opposite faces exert a net force, and its acceleration is
inversely proportional to the mass density.  The result has the form
\begin{equation}
 c_s^2=
 \left(\frac{\partial P}{\partial \rho_m}\right)_{\cal C},
 \qquad
 \rho_m=m\rho,
 \label{eq:sm7-sound-general}
\end{equation}
where \(\rho\) is number density and the subscript \({\cal C}\) specifies the
thermodynamic condition followed by the disturbance.  The board calculation
uses
\begin{equation}
 P=\rho k_B T
 \quad\Longrightarrow\quad
 \frac{1}{m}
 \left(\frac{\partial P}{\partial\rho}\right)_T
 =\frac{k_B T}{m}.
 \label{eq:sm7-sound-isothermal}
\end{equation}
This is the isothermal sound-speed estimate.  It differs from
\eqref{eq:sm7-molecular-speed} by a factor of \(3\) in the squared speed, or
\(\sqrt{3}\) in the speed.

\begin{figure}[t]
\centering
\includegraphics[width=0.90\textwidth]{lecture_07_figure_02.png}
\caption{The molecular-speed estimate, the pressure derivative, and the
ideal-gas substitution displayed together.}
\lectureframe{00:08:55}
\end{figure}

There is an important refinement.  A small-amplitude sound wave normally
compresses and rarefies the gas faster than heat can flow between neighboring
elements.  The local motion is therefore approximately adiabatic, not
isothermal:
\begin{equation}
 c_s^2=
 \left(\frac{\partial P}{\partial\rho_m}\right)_S.
 \label{eq:sm7-sound-adiabatic}
\end{equation}
For an ideal gas with adiabatic index \(\gamma\),
\begin{equation}
 P\rho_m^{-\gamma}=\text{constant},
 \qquad
 c_s^2=\gamma\frac{P}{\rho_m}
       =\gamma\frac{k_B T}{m}.
 \label{eq:sm7-sound-ideal-adiabatic}
\end{equation}
For room-temperature diatomic air, \(\gamma\simeq 7/5\).  Thus the
fixed-temperature calculation is a useful first estimate, while the
adiabatic derivative supplies the physically appropriate refinement.

Using the rough values
\begin{equation}
\begin{aligned}
 k_B&\simeq1.3\times10^{-23}\ {\rm J\,K^{-1}},
 &T&\simeq300\ {\rm K},\\
 m_{\rm air}&\simeq4\times10^{-26}\ {\rm kg}.&&
\end{aligned}
 \label{eq:sm7-air-data}
\end{equation}
the molecular estimate is about \(500\ {\rm m\,s^{-1}}\).  Temperature here
must be an absolute temperature in kelvins, not an offset scale such as
Fahrenheit.  The familiar sound
speed is of order \(300\ {\rm m\,s^{-1}}\), and the adiabatic ideal-gas
formula gives about \(340\ {\rm m\,s^{-1}}\) at room temperature.  The rough
calculation has therefore found the correct physical scale.

\begin{figure}[t]
\centering
\includegraphics[width=0.90\textwidth]{lecture_07_figure_03.png}
\caption{The room-air estimate using Boltzmann's constant, an effective
molecular mass near \(30\) proton masses, and \(T\simeq300\ {\rm K}\).}
\lectureframe{00:10:55}
\end{figure}

Long before a microscopic account was available, one could determine the
sound speed from measurements of pressure as a function of density.  The
statistical argument adds the striking interpretation that the resulting
macroscopic speed is set by the thermal speed of individual molecules.

\begin{classroomqa}
\audiencequestion{Should neutrons be added to the molecular mass, and does
the number \(14\) for nitrogen count only protons?}
\lecturerresponse{The quoted atomic weight already counts protons and
neutrons.  Molecular nitrogen contains two atoms of atomic weight about
\(14\), so its mass is about \(28\) proton masses; rounding \(28\) to \(30\)
is sufficient for this estimate.
\lecturetimestamp{00:10:58--00:11:16}}
\end{classroomqa}

\begin{classroomqa}
\audiencequestion{The factor \(3\) came from three translational directions,
whereas sound propagates along one direction.  Should that remove the factor
of \(3\)?}
\lecturerresponse{It correctly warns us not to identify the three-dimensional
root-mean-square molecular speed with a one-dimensional propagation speed.
The opening argument was only an order-of-magnitude picture.  A controlled
answer comes from the pressure derivative, together with the thermodynamic
condition appropriate to the wave.
\lecturetimestamp{00:11:12--00:12:00}}
\end{classroomqa}

\begin{classroomqa}
\audiencequestion{Did the derivative assume that temperature remains
constant during the wave?}
\lecturerresponse{The displayed ideal-gas substitution did.  Holding
temperature fixed is a plausible provisional estimate for a very small
disturbance; the standard refinement is that ordinary sound is approximately
adiabatic, which replaces the isothermal derivative by
\(\left(\partial P/\partial\rho_m\right)_S\).
\lecturetimestamp{00:12:59--00:14:01}}
\end{classroomqa}

\begin{classroomqa}
\audiencequestion{If the pressure and density variations are both small,
does that make their derivative small?}
\lecturerresponse{No.  A derivative is the ratio of two small variations in
the infinitesimal limit.  Both \(\delta P\) and \(\delta\rho\) may be small
while \(\delta P/\delta\rho\) remains finite and sets the wave speed.
\lecturetimestamp{00:14:03--00:14:23}}
\end{classroomqa}

\begin{classroomqa}
\audiencequestion{Should the speed of sound rise when the gas is compressed,
as it does when one passes from a gas toward a liquid or solid?}
\lecturerresponse{Yes.  The derivative is evaluated at the state under
consideration, and \(P(\rho,T)\) need not remain linear.  When molecules are
squeezed close together, the equation of state becomes nonideal and
\(\partial P/\partial\rho\) can grow substantially.
\lecturetimestamp{00:14:25--00:15:13}}
\end{classroomqa}

\begin{classroomqa}
\audiencequestion{Can the value \(500\ {\rm m\,s^{-1}}\) be used at every
pressure?}
\lecturerresponse{No.  It belongs to the low-density range in which the
ideal-gas equation is a useful approximation.  Once higher powers of density
enter \(P(\rho,T)\), the derivative and therefore the sound speed depend on
the pressure and density.
\lecturetimestamp{00:15:14--00:16:29}}
\end{classroomqa}

\begin{classroomqa}
\audiencequestion{Does the original picture of sound spreading at roughly
the molecular speed also fail as density increases?}
\lecturerresponse{Probably.  That picture is tailored to a dilute gas.
Collective interactions become increasingly important at high density, which
is precisely why the equation-of-state derivative is the more general
description.
\lecturetimestamp{00:16:32--00:16:53}}
\end{classroomqa}

\begin{classroomqa}
\audiencequestion{Why was the nitrogen-molecule mass written as roughly
\(30\) proton masses rather than \(28\)?}
\lecturerresponse{Only numerical accuracy appropriate to the estimate is
needed, so \(28\) was rounded to \(30\).  Dividing the original
\(500\ {\rm m\,s^{-1}}\) estimate by \(\sqrt{3}\simeq1.7\) also brings it
close to \(300\ {\rm m\,s^{-1}}\), although the adiabatic derivative is the
controlled way to calculate the correction.
\lecturetimestamp{00:16:54--00:17:27}}
\end{classroomqa}

\section{One Classical Oscillator in a Heat Bath}

We now replace the freely translating molecule by one harmonic oscillator and
treat everything else as a heat bath.  A mass on a spring is the immediate
picture, but the model is much broader.  An electromagnetic cavity mode
oscillates; atoms in a crystal oscillate; almost any stable system displaced
slightly from equilibrium executes harmonic motion to leading order.

With \(x\) measured from equilibrium, momentum \(p\), mass \(m\), and spring
constant \(\kappa\), the Hamiltonian is
\begin{equation}
 H(x,p)=\frac{p^2}{2m}+\frac{\kappa x^2}{2},
 \qquad
 \omega=\sqrt{\frac{\kappa}{m}}.
 \label{eq:sm7-classical-hamiltonian}
\end{equation}
The bath fixes \(T\), and from this point we again use \(k_B=1\), so
\(\beta=1/T\).  An oscillator equilibrated with the room has a temperature
near \(300\ {\rm K}\); one immersed in liquid helium has a temperature of only
a few kelvins.  A putative spring placed in the center of the Sun would not
survive as a spring, but the example makes the role of the environment clear.
The Boltzmann weight factorizes:
\begin{equation}
 e^{-\beta H}
 =e^{-\beta p^2/(2m)}
  e^{-\beta\kappa x^2/2}.
 \label{eq:sm7-factorized-weight}
\end{equation}

\begin{figure}[t]
\centering
\includegraphics[width=0.90\textwidth]{lecture_07_figure_04.png}
\caption{A spring-mass oscillator and its kinetic-plus-potential Hamiltonian.}
\lectureframe{00:20:55}
\end{figure}

Unlike the free gas, the position integral does not merely yield the
container volume.  The restoring potential weights each displacement:
\begin{equation}
 Z_{\rm cl}
 =\frac{1}{h}\int_{-\infty}^{\infty}dp
             \int_{-\infty}^{\infty}dx\,
 e^{-\beta p^2/(2m)}
 e^{-\beta\kappa x^2/2}.
 \label{eq:sm7-classical-z}
\end{equation}
The factor \(1/h\) makes a one-dimensional classical partition function
dimensionless.  It is suppressed in the board calculation because it is
independent of \(\beta\) and cannot affect the mean energy.

\begin{figure}[t]
\centering
\includegraphics[width=0.90\textwidth]{lecture_07_figure_05.png}
\caption{The classical oscillator partition function separated into its
momentum and displacement Gaussian integrals.}
\lectureframe{00:22:55}
\end{figure}

For the momentum integral, set
\begin{equation}
 q=\sqrt{\frac{\beta}{2m}}\,p,
 \qquad
 dp=\sqrt{\frac{2m}{\beta}}\,dq.
 \label{eq:sm7-p-gaussian-substitution}
\end{equation}
For the position integral, set
\begin{equation}
 y=\sqrt{\frac{\beta\kappa}{2}}\,x,
 \qquad
 dx=\sqrt{\frac{2}{\beta\kappa}}\,dy.
 \label{eq:sm7-x-gaussian-substitution}
\end{equation}
Since \(\int_{-\infty}^{\infty}e^{-u^2}\,du=\sqrt{\pi}\), the
unnormalized phase-space integral used on the board is
\begin{equation}
 \int dp\,dx\,e^{-\beta H}
 =\frac{2\pi}{\beta}\sqrt{\frac{m}{\kappa}}
 =\frac{2\pi}{\beta\omega}.
 \label{eq:sm7-classical-z-unnormalized}
\end{equation}
With the phase-space cell restored,
\begin{equation}
 Z_{\rm cl}=\frac{2\pi}{h\beta\omega}
           =\frac{1}{\beta\hbar\omega}.
 \label{eq:sm7-classical-z-normalized}
\end{equation}

Only the \(\beta\)-dependence matters for the energy:
\begin{equation}
 \begin{aligned}
 \log Z_{\rm cl}&=\text{constant}-\log\beta,\\
 \expect{H}
 &=-\frac{\partial\log Z_{\rm cl}}{\partial\beta}\\
 &=\frac{1}{\beta}=T.
 \end{aligned}
 \label{eq:sm7-classical-energy}
\end{equation}

\begin{figure}[t]
\centering
\includegraphics[width=0.90\textwidth]{lecture_07_figure_06.png}
\caption{The complete classical chain on the right board: Gaussian
factorization, \(Z\propto(\beta\omega)^{-1}\), and
\(\expect{H}=1/\beta=T\).}
\lectureframe{00:30:15}
\end{figure}

There is no factor of \(3\) because this is one spatial oscillator.  There is
also no remaining factor of \(1/2\), because there are two independent
quadratic terms.  Each Gaussian contributes \(T/2\):
\begin{equation}
 \expect{K}=\frac{T}{2},
 \qquad
 \expect{U}=\frac{T}{2},
 \qquad
 \expect{H}=T.
 \label{eq:sm7-equipartition-oscillator}
\end{equation}
This is equipartition in its most useful form: every independent quadratic
term in the classical Hamiltonian contributes \(T/2\).

The answer is independent of both \(m\) and \(\kappa\).  The motion itself
does depend on those parameters:
\begin{equation}
 \expect{p^2}=mT,
 \qquad
 \expect{x^2}=\frac{T}{\kappa},
 \qquad
 \expect{\dot{x}^2}=\frac{T}{m}.
 \label{eq:sm7-classical-fluctuations}
\end{equation}
Thus a heavier mass moves more slowly and a stiffer spring has a smaller
amplitude, even though the mean energy remains \(T\).
That becomes disturbing when \(\kappa\) is made fantastically large.  A
spring too stiff to move ought to look like a frozen constraint, yet
classical equipartition still assigns it energy \(T\).

\begin{classroomqa}
\audiencequestion{As \(\kappa\to\infty\), the partition function tends to
zero.  How can its logarithm still be used?}
\lecturerresponse{Keep \(\kappa\) finite but arbitrarily large.  Its
contribution to \(\log Z\) is additive and independent of \(\beta\), so the
\(\beta\)-derivative still gives \(T\).  The limiting discomfort is real:
the classical calculation is internally consistent, but it omits a feature
of nature needed for very stiff oscillators.
\lecturetimestamp{00:34:01--00:34:54}}
\end{classroomqa}

\begin{classroomqa}
\audiencequestion{The ideal-gas partition function required an \(N!\).
Should there be an analogous factorial here?}
\lecturerresponse{No.  We are studying one specified oscillator, not
permuting \(N\) identical particles.  For many independent labeled modes the
partition functions multiply.  An indistinguishability factor belongs to a
different counting problem.
\lecturetimestamp{00:35:03--00:35:40}}
\end{classroomqa}

\section{Why Classical Equipartition Was a Crisis}

Real molecules are not structureless points.  A diatomic molecule contains a
bond that can vibrate, hence an internal oscillator.  A naive classical count
would assign that mode additional thermal energy even at ordinary
temperatures.  In a deliberately schematic count, activating one vibrational
oscillator raises an energy of \(3T/2\) toward \(5T/2\).  Precise
heat-capacity bookkeeping must also include rotational degrees of freedom and
distinguish energy from its temperature derivative.  The decisive
experimental fact is simpler:
stiff internal modes did not contribute at low temperature as classical
equipartition demanded.

One could deny that the molecule has internal structure, which is absurd, or
recognize that classical mechanics is missing an ingredient.  The ingredient
is quantum mechanics.  We need very little of it: only that a harmonic
oscillator has equally spaced energy levels.

\section{The Quantum Oscillator}

First shift the energy origin and write the excitation spectrum as
\begin{equation}
 E_n^{\rm exc}=n\hbar\omega,
 \qquad n=0,1,2,\ldots.
 \label{eq:sm7-shifted-spectrum}
\end{equation}
The partition function is now a genuine sum rather than a phase-space
integral:
\begin{align}
 Z_{\rm exc}
 &=\sum_{n=0}^{\infty}e^{-\beta n\hbar\omega} \notag\\
 &=1+x+x^2+\cdots
  =\frac{1}{1-x},
 \qquad x=e^{-\beta\hbar\omega}, \notag\\
 &=\frac{1}{1-e^{-\beta\hbar\omega}}.
 \label{eq:sm7-quantum-z}
\end{align}

\begin{figure}[t]
\centering
\includegraphics[width=0.90\textwidth]{lecture_07_figure_07.png}
\caption{Equally spaced levels and the geometric-series evaluation of the
quantum oscillator partition function.}
\lectureframe{00:42:05}
\end{figure}

Differentiating gives the mean excitation energy:
\begin{align}
 E_{\rm exc}
 &=-\frac{\partial\log Z_{\rm exc}}{\partial\beta}
  =-\frac{1}{Z_{\rm exc}}
    \frac{\partial Z_{\rm exc}}{\partial\beta}
 \notag\\
 &=\hbar\omega\,
   \frac{e^{-\beta\hbar\omega}}
        {1-e^{-\beta\hbar\omega}}
  =\frac{\hbar\omega}{e^{\beta\hbar\omega}-1}.
 \label{eq:sm7-quantum-energy}
\end{align}

\begin{figure}[t]
\centering
\includegraphics[width=0.90\textwidth]{lecture_07_figure_08.png}
\caption{The thermal excitation energy obtained from the
\(\beta\)-derivative of the quantum partition function.}
\lectureframe{00:46:05}
\end{figure}

\subsection{High Temperature}

When \(\beta\hbar\omega\ll1\),
\begin{equation}
 e^{-\beta\hbar\omega}
 =1-\beta\hbar\omega
  +O\!\left((\beta\hbar\omega)^2\right).
 \label{eq:sm7-high-t-expansion}
\end{equation}
The numerator of \eqref{eq:sm7-quantum-energy} approaches
\(\hbar\omega\), while the denominator is
\(\beta\hbar\omega\) to leading order.  Therefore
\begin{equation}
 E_{\rm exc}\simeq\frac{1}{\beta}=T.
 \label{eq:sm7-high-t-energy}
\end{equation}
The quantum answer recovers classical equipartition when many level spacings
fit inside the thermal energy.

\subsection{Low Temperature}

When \(\beta\hbar\omega\gg1\), the exponential is tiny:
\begin{equation}
 E_{\rm exc}
 \simeq\hbar\omega e^{-\beta\hbar\omega}.
 \label{eq:sm7-low-t-energy}
\end{equation}
The oscillator's thermal excitation is exponentially suppressed.  In the
shifted convention it is almost zero; the physical zero-point energy will be
restored below.

\begin{figure}[t]
\centering
\includegraphics[width=0.90\textwidth]{lecture_07_figure_09.png}
\caption{The low-temperature expansion: the first thermal excitation is
suppressed by \(e^{-\beta\hbar\omega}\).}
\lectureframe{00:51:55}
\end{figure}

The crossover occurs when
\begin{equation}
 \beta\hbar\omega\sim1,
 \qquad\text{or equivalently}\qquad
 T\sim\hbar\omega.
 \label{eq:sm7-crossover}
\end{equation}
This condition has a direct meaning.  Classical equipartition would give the
oscillator energy \(T\).  Quantum effects become unavoidable when that scale
falls below one level spacing.  Thus
\begin{equation}
\begin{aligned}
 \hbar\omega\gg T&:\quad \text{frozen and quantum},\\
 \hbar\omega\ll T&:\quad \text{active and classical}.
\end{aligned}
 \label{eq:sm7-regimes}
\end{equation}
A stiffer spring has larger \(\omega=\sqrt{\kappa/m}\), so it requires a
higher temperature before it behaves classically.

\section{Degrees of Freedom Appear as We Heat Matter}

At low temperature, a diatomic molecule can thermally resemble a point
particle because its stiff bond vibration is frozen.  As \(T\) rises through
\(\hbar\omega\), the vibrational mode activates.  Rotational modes may
activate at a lower scale; still higher temperatures reveal more internal
motion and eventually ionize the atom.  Heating does not merely put more
energy into a fixed list of modes.  It reveals degrees of freedom that were
invisible at lower temperature.

An imaginary nested molecule sharpens the point.  Two atoms are
joined by a stiff spring, while one apparent atom is itself a pair joined by
an even stiffer spring.  At low \(T\), both structures look pointlike.  At the
first threshold the outer vibration appears; at the second threshold the
inner structure appears.  The complexity of the system is uncovered one
energy scale at a time.

\begin{classroomqa}
\audiencequestion{Once an oscillator is already in the classical regime,
does further heating tighten the spring or change its frequency?}
\lecturerresponse{No.  The spring constant and frequency remain fixed.
Higher temperature produces larger typical momenta and amplitudes, hence
larger energy.  Stiffness determines the temperature at which the classical
law begins, not a temperature-dependent tightening of the spring.
\lecturetimestamp{01:03:17--01:05:11}}
\end{classroomqa}

\subsection{Crystals and Einstein's Specific Heat}

A crystal contains a great many vibrational modes.  Classical equipartition
assigns \(T\)-scale energy to every active oscillator and therefore
overestimates the low-temperature heat capacity.  Einstein's 1907 model
treated crystal vibrations quantum mechanically and showed why they are
suppressed until thermal energy reaches their level spacing.  More accurate
solid-state models allow a distribution of mode frequencies, but the
freeze-out principle is already present in Einstein's calculation.

\begin{classroomqa}
\audiencequestion{Was this the work for which Einstein received the Nobel
Prize, and did it come only after quantum mechanics was established?}
\lecturerresponse{No.  The Nobel citation concerned the photoelectric
effect.  The specific-heat model was published in 1907, long before the full
foundations of quantum mechanics were understood.  Einstein accepted and
developed crucial quantum ideas even though he later objected to aspects of
the mature theory.
\lecturetimestamp{00:58:38--01:03:16}}
\end{classroomqa}

\begin{classroomqa}
\audiencequestion{Classical statistical mechanics replaced a sum over
discrete states by an apparently uniform integral over position and
momentum.  What justifies that state counting?}
\lecturerresponse{The ultimate justification comes from quantum mechanics.
In the high-temperature limit,
\[
 Z_{\rm exc}\simeq\frac{1}{\beta\hbar\omega},
\]
which exactly matches the dimensionless classical result
\eqref{eq:sm7-classical-z-normalized}.  Quantum cells in phase space become
the classical phase-space measure when the level spacing is unresolved.
\lecturetimestamp{00:59:38--01:03:06}}
\end{classroomqa}

\subsection{A String Contains Infinitely Many Oscillators}

A violin or guitar string decomposes into normal modes: the fundamental,
higher harmonics, and progressively shorter wavelengths.  Mathematically
there are infinitely many oscillator modes.  If classical equipartition
assigned energy \(T\) to every one, the string in thermal equilibrium would
have infinite energy.  Quantum mechanics prevents that catastrophe.  The
shortest-wavelength modes have the highest frequencies and remain frozen at
any finite temperature; progressively more modes activate as \(T\) rises.
The same logic will later underlie blackbody radiation.

\begin{figure}[t]
\centering
\includegraphics[width=0.90\textwidth]{lecture_07_figure_10.png}
\caption{Several normal modes superposed on a string, illustrating why a
continuous string is an infinite collection of harmonic oscillators.}
\lectureframe{01:09:30}
\end{figure}

\begin{classroomqa}
\audiencequestion{Can the quantum-to-classical crossover be observed
directly, and can heating reverse the low-temperature behavior?}
\lecturerresponse{Yes.  The temperature dependence of a crystal's specific
heat makes the freeze-out and subsequent activation of vibrational modes
plainly measurable.
\lecturetimestamp{01:09:11--01:10:13}}
\end{classroomqa}

\begin{classroomqa}
\audiencequestion{For many oscillators, do we repeat the single-oscillator
partition function for each mode?}
\lecturerresponse{For independent modes,
\[
\begin{aligned}
 Z_{\rm tot}&=\prod_j Z_j,&
 \log Z_{\rm tot}&=\sum_j\log Z_j,\\
 E_{\rm tot}&=\sum_j E_j.&&
\end{aligned}
\]
Modes with \(\hbar\omega_j\gg T\) make exponentially small thermal
contributions; those below the threshold contribute appreciably.
\lecturetimestamp{01:10:14--01:11:00}}
\end{classroomqa}

\begin{classroomqa}
\audiencequestion{What happened to the oscillator's zero-point energy
\(\hbar\omega/2\)?}
\lecturerresponse{The physical levels are
\[
 E_n=\left(n+\frac12\right)\hbar\omega.
\]
Consequently
\[
 \begin{aligned}
 Z_{\rm phys}
 &=e^{-\beta\hbar\omega/2}Z_{\rm exc},\\
 E_{\rm phys}
 &=\frac{\hbar\omega}{2}
  +\frac{\hbar\omega}{e^{\beta\hbar\omega}-1}.
 \end{aligned}
\]
The common shift does not change level probabilities, entropy, heat
capacity, or thermal excitation energy, but it does shift the absolute mean
energy by \(\hbar\omega/2\).  The multiplying factor in \(Z\) is
\(\beta\)-dependent, so it must not be discarded when absolute energy is
wanted.
\lecturetimestamp{01:11:08--01:11:37}}
\end{classroomqa}

\section{The Second Law as a Reversibility Puzzle}

We can now return to the promised problem.  The second law says that ordinary
macroscopic evolution is irreversible and that entropy tends to increase.
Microscopic Hamiltonian evolution is reversible.  How can a reversible map
of exact states produce an irreversible law of macroscopic states?

Represent the system by a point in a high-dimensional phase space,
\[
 \Gamma=(q_1,\ldots,q_N,p_1,\ldots,p_N).
\]
Our knowledge is not usually an exact point.  Suppose it is a uniform
probability density on a compact patch \(A\) of volume \(\Omega_A\):
\begin{equation}
 f(\Gamma)=
 \begin{cases}
  1/\Omega_A,&\Gamma\in A,\\
  0,&\Gamma\notin A.
 \end{cases}
 \label{eq:sm7-uniform-blob}
\end{equation}
In units with \(k_B=1\), its fine-grained Gibbs entropy is
\begin{equation}
 S_{\rm fine}
 =-\int d\Gamma\,f\log f
 =\log\Omega_A.
 \label{eq:sm7-blob-entropy}
\end{equation}
A smaller patch means more information and lower entropy.  This fixes the
sign: increasing uncertainty enlarges the occupied region and raises
\(\log\Omega_A\).

\begin{figure}[t]
\centering
\includegraphics[width=0.90\textwidth]{lecture_07_figure_11.png}
\caption{A compact probability patch in a schematic two-dimensional slice
of the full position-momentum phase space.}
\lectureframe{01:20:20}
\end{figure}

\begin{classroomqa}
\audiencequestion{Is the system physically constrained to remain inside the
drawn patch, and is it already in equilibrium?}
\lecturerresponse{No.  The patch describes what is known about the initial
state, not a wall in phase space.  It need not be an equilibrium
distribution.  A gas initially confined to a tank in one corner of a room is
a concrete example of special, nonequilibrium information.
\lecturetimestamp{01:14:55--01:15:24}}
\end{classroomqa}

\section{Liouville Flow and Coarse-Grained Growth}

Liouville's theorem says that Hamiltonian flow preserves phase-space volume.
The exact patch can translate, stretch, fold, and become extremely thin, but
\begin{equation}
 \Omega_{A(t)}=\Omega_{A(0)},
 \qquad
 S_{\rm fine}(t)=S_{\rm fine}(0).
 \label{eq:sm7-liouville}
\end{equation}
The discrete analog is a reversible permutation of states.  If one initially
knows that a die is in one of two states, a one-to-one evolution may change
which two states they are, but it does not change the number of possibilities.

Real descriptions have finite resolution.  Quantum mechanics supplies a
natural phase-space cell, while a classical experiment has finite measuring
power even without invoking quantum theory.  Under Hamiltonian flow the patch
can become a long, thin snake.  If its separation between neighboring folds
falls below the resolution scale, we replace each unresolved neighborhood by
a coarse cell.  The exact occupied volume remains fixed, but the union of
coarse cells grows:
\begin{equation}
 \Omega_{\rm coarse}(t)\geq\Omega_{A(t)},
 \qquad
 S_{\rm coarse}(t)=\log\Omega_{\rm coarse}(t).
 \label{eq:sm7-coarse-entropy}
\end{equation}

\begin{figure}[t]
\centering
\includegraphics[width=0.90\textwidth]{lecture_07_figure_12.png}
\caption{A volume-preserving filament viewed through finite-resolution
cells.  The exact black filament keeps its area; the red coarse cells occupy
a larger effective region.}
\lectureframe{01:22:10}
\end{figure}

\begin{classroomqa}
\audiencequestion{The coarse cells overlap.  Does that violate Liouville's
theorem?}
\lecturerresponse{No.  Liouville governs the exact filament, which remains
one-to-one and volume preserving.  Overlap belongs to our blurred
description: we have discarded information about which thin strand contains
the state.
\lecturetimestamp{01:21:32--01:21:50}}
\end{classroomqa}

A useful analogy is cotton.  The material volume of all fibers can
remain nearly fixed while a teased-out wad occupies a large apparent volume
to an observer unable to resolve individual fibers.  Likewise, a
phase-space patch grows finer tentacles; those tentacles develop still finer
tentacles.  Any fixed finite resolution eventually fails to follow them.

\begin{classroomqa}
\audiencequestion{Does entropy merely appear to increase classically, while
quantum mechanics makes the increase real by imposing a fundamental
resolution limit?}
\lecturerresponse{The coarse-graining argument does not require a quantum
limit.  Entropy refers jointly to the system and the macroscopic distinctions
retained in its description.  With infinite microscopic information the
fine-grained entropy is constant; with any fixed finite resolution, chaotic
filaments eventually become unresolved and the coarse-grained entropy
increases.
\lecturetimestamp{01:24:04--01:25:20}}
\end{classroomqa}

\section{Chaos, a Finite Accessible Region, and Equilibration}

For a gas in a box, positions are bounded.  At fixed finite total energy,
momenta are bounded as well, even if all energy momentarily concentrates in
one particle.  The dynamically accessible energy shell therefore has finite
measure under the usual confinement assumptions.

\begin{figure}[t]
\centering
\includegraphics[width=0.90\textwidth]{lecture_07_figure_13.png}
\caption{A bounded accessible phase-space region beside the filament and
its coarse-grained representation.}
\lectureframe{01:28:10}
\end{figure}

Nearby initial points in a chaotic system follow similar histories for a
while and then separate.  In a billiard break, an imperceptible change in the
cue-ball direction is copied into later collisions; the differences compound
until the final configurations are unrelated.  A whole initial patch is
therefore stretched and folded through the accessible region.  At finite
resolution it approaches a broad, nearly uniform equilibrium description.

\begin{figure}[t]
\centering
\includegraphics[width=0.90\textwidth]{lecture_07_figure_14.png}
\caption{A compact special state evolved into a complicated filament
through the bounded accessible region while the coarse-grained cells spread
over a much larger area.}
\lectureframe{01:33:10}
\end{figure}

Two qualifications matter.  Chaos by itself does not prove that every
trajectory visits every accessible neighborhood; that stronger sampling
claim requires an ergodic or mixing property.  Nor is every physical system
chaotic.  A generic many-body gas and idealized billiards provide the intended
setting, where mixing is a useful physical expectation.

\section{Recurrence Does Not Abolish Irreversibility}

Suppose all air molecules begin in one corner with a specially prepared
velocity pattern.  Collisions quickly spread them through the room, and the
coarse-grained description reaches equilibrium.  In a finite
measure-preserving phase space, the Poincare recurrence theorem says that
almost every initial point returns arbitrarily close to its starting point.
For a tolerance \(\epsilon>0\), a recurrence time \(t_R(\epsilon)\) is finite
under the theorem's hypotheses, though for a macroscopic system it is
fantastically large.  Exact return to an infinitely precise point is neither
required nor generally expected.

\begin{classroomqa}
\audiencequestion{In the billiard example, does striking the cue ball inject
energy and spoil the closed evolution?}
\lecturerresponse{The chosen initial state begins immediately after the
strike, with that energy already included.  Recurrence concerns the
subsequent isolated idealized billiard dynamics; a return near the initial
state includes a configuration equivalent to the cue ball approaching the
rack within the chosen tolerance.
\lecturetimestamp{01:35:37--01:35:55}}
\end{classroomqa}

\begin{classroomqa}
\audiencequestion{If the gas can eventually return to the corner, have we
lost information, and is the system really reversible?}
\lecturerresponse{The microscopic dynamics remains reversible throughout.
Coarse-grained entropy can rise because we stop resolving fine correlations,
yet an exceptionally rare fluctuation can return the system near its special
initial macrostate.  Recurrence and coarse-grained irreversibility are
therefore compatible.
\lecturetimestamp{01:37:34--01:38:15}}
\end{classroomqa}

This leads to Boltzmann's essential formulation.  Entropy does not increase
at every logically possible instant.  Given a nonequilibrium macrostate of a
large system, overwhelmingly more compatible microscopic continuations lead
toward larger macrostates than toward still more special ones.  Entropy
\emph{probably} increases.  Downward fluctuations and recurrences are
allowed, but their probabilities and time scales make them invisible in
ordinary macroscopic experience.  The reversibility problem was one of the
great historical struggles of statistical mechanics and weighed heavily on
Boltzmann.  That historical connection should not be turned into a simple
causal claim about his life.

\section{What Chaos Means}

For a more operational definition of chaos, suppose
two initial phase points differ by \(\delta\Gamma(0)\), a chaotic direction
typically has
\begin{equation}
 \lVert\delta\Gamma(t)\rVert
 \sim e^{\lambda t}\lVert\delta\Gamma(0)\rVert,
 \qquad \lambda>0,
 \label{eq:sm7-lyapunov-growth}
\end{equation}
over the regime where linearized separation is meaningful.  The number
\(\lambda\) is a Lyapunov exponent.  To predict a time \(t\), the required
precision in initial data scales roughly as \(e^{-\lambda t}\).  Precision in
the dynamical law matters too: model errors are amplified along with errors
in initial conditions.  Weather is the familiar example.

\begin{classroomqa}
\audiencequestion{Is there a principle saying that \emph{most} systems are
chaotic, or is that merely empirical?}
\lecturerresponse{There are infinitely many chaotic and nonchaotic
Hamiltonians, and even the word \emph{most} requires a measure on the space
of Hamiltonians.  Complicated interacting systems are often chaotic, but
there is no elementary universal test that classifies every Hamiltonian by
inspection.
\lecturetimestamp{01:39:27--01:40:49}}
\end{classroomqa}

\begin{classroomqa}
\audiencequestion{What precisely does \emph{chaotic} mean in this
discussion?}
\lecturerresponse{Nearby phase trajectories separate rapidly enough that
long-time prediction requires exponentially improving knowledge of initial
conditions.  A harmonic oscillator is the contrast: neighboring phase points
rotate together and remain close.  The exact two-body Kepler problem is also
nonchaotic, while the general three-body problem has chaotic regimes.
\lecturetimestamp{01:40:50--01:43:45}}
\end{classroomqa}

\begin{classroomqa}
\audiencequestion{Must the laws of motion, as well as the initial
conditions, be known to high precision?}
\lecturerresponse{Yes.  Both sources of error propagate through the
dynamics, and chaotic evolution amplifies both.
\lecturetimestamp{01:43:46--01:43:55}}
\end{classroomqa}

The double pendulum supplies a vivid example.  An ordinary pendulum is
integrable and predictable.  Add a second pendulum, and trajectories that
begin almost together can eventually transfer energy differently and head
in entirely different directions.

\begin{figure}[t]
\centering
\includegraphics[width=0.90\textwidth]{lecture_07_figure_15.png}
\caption{The double pendulum as a concrete example of deterministic chaos.}
\lectureframe{01:44:40}
\end{figure}

\begin{classroomqa}
\audiencequestion{Why can one mechanical system be chaotic while another is
not?}
\lecturerresponse{That is a property of the Hamiltonian and its invariant
structures, not a consequence of friction or experimental noise.  Determining
it can be mathematically difficult; complexity is suggestive but not a
proof.
\lecturetimestamp{01:44:03--01:45:58}}
\end{classroomqa}

\begin{classroomqa}
\audiencequestion{Is an ideal double pendulum still chaotic with no
friction?}
\lecturerresponse{Yes.  Dissipation is not required.  Exact deterministic
Hamiltonian motion can have sensitive dependence on initial conditions.
Infinite-precision initial data would determine the trajectory indefinitely,
but finite precision gives a finite prediction horizon.
\lecturetimestamp{01:45:57--01:47:54}}
\end{classroomqa}

\begin{classroomqa}
\audiencequestion{Is chaos the same as having no closed-form solution?}
\lecturerresponse{No.  Chaotic systems generally lack a globally useful
closed-form trajectory, but analytic solvability and sensitive dependence are
different properties.  Chaos is diagnosed dynamically by trajectory
separation, not merely by our inability to write a formula.
\lecturetimestamp{01:48:20--01:49:23}}
\end{classroomqa}

\begin{classroomqa}
\audiencequestion{How is pervasive chaos different from balancing exactly
at the unstable top of a hill?}
\lecturerresponse{The hilltop is one exceptional unstable state: away from
it, nearby trajectories need not separate rapidly.  In a chaotic region,
instability is typical along the orbit and nearby phase points repeatedly
separate.  A positive Lyapunov exponent captures that persistent growth.
\lecturetimestamp{01:49:27--01:50:11}}
\end{classroomqa}

\section{From Microscopic Motion to Statistical Law}

The speed-of-sound estimate showed how microscopic thermal motion fixes a
macroscopic propagation scale, while the questions exposed the
thermodynamic derivative that must be specified.  The oscillator then showed
the reach and the failure of classical equipartition.  Quantum level spacing
freezes stiff modes, explains why heat capacities change with temperature,
and prevents infinitely many high-frequency modes from carrying infinite
thermal energy.

The same progression from exact microphysics to effective macroscopic
description resolves the apparent contradiction at the heart of the second
law.  Hamiltonian flow preserves exact phase-space volume and remains
reversible.  Chaotic stretching creates finer correlations than a
macroscopic description retains, so the coarse-grained occupied region
grows.  Recurrence guarantees that rare returns are not forbidden; probability
explains why they do not define ordinary experience.  The second law is not a
new irreversible microscopic force.  It is the overwhelmingly reliable
statistical direction of evolution for special macroscopic initial
conditions in systems with enormously many degrees of freedom.
